% GENERATED FILE. Edit canonical lessons, then run npm run generate:books.
% canonical-source-hash: fnv1a64:61e5d87f2d76eebf
% canonical-lessons: BN-C162-sugondhi, BN-C162-pocha, BN-C162-adha, BN-C162-puro, BN-C162-samanyo

\chapter{Fragrant, Rotten, Half, Whole, Slight}
\label{ch:bn-fragrant-rotten-half-whole-slight}
\hlchaptermodality{162}

\begin{chapteropening}
\textbf{By the end of this chapter:} \emph{I can say fragrant, rotten, half, whole and slight.}

Complete the last lesson of chapter 162: I can say fragrant, rotten, half, whole and slight.
\end{chapteropening}

\section[\texorpdfstring{\bn{সুগন্ধি} (sugôndhi)}{সুগন্ধি (sugôndhi)}]{\bn{সুগন্ধি} (sugôndhi) — fragrant}

\label{lesson:BN-C162-sugondhi}



\begin{quote}
Before the new one: say the Bengali for contented, then the Bengali for annoyed.
\end{quote}

\subsection*{You'll want to know: \bn{সুগন্ধি}}
\textbf{\bn{সুগন্ধি}} — \emph{sugôndhi} — \textquotedblleft{}fragrant\textquotedblright{}.

\begin{grammarlens}[title={What you've built}]
One more word for this chapter.
\end{grammarlens}

\subsection*{Your turn}
\begin{itemize}
\raggedright
  \item \emph{Say it:} \emph{sugôndhi}
  \item \emph{Say it:} \emph{sugôndhi}, once more
  \item \emph{Say it:} say \emph{sugôndhi}
  \item \emph{Recall:} say the Bengali for contented, then the Bengali for annoyed, then say \emph{sugôndhi} again
\end{itemize}

\begin{tcolorbox}[breakable,colback=teal!4,colframe=teal!35!black,title={Before you move on}]
What does \bn{সুগন্ধি} mean? (\textquotedblleft{}Fragrant\textquotedblright{}.) Say it once more.
\end{tcolorbox}
\section[\texorpdfstring{\bn{পচা} (pôchā)}{পচা (pôchā)}]{\bn{পচা} (pôchā) — rotten}

\label{lesson:BN-C162-pocha}



\begin{quote}
Before the new one: say the Bengali for annoyed, then the Bengali for fragrant.
\end{quote}

\subsection*{You'll want to know: \bn{পচা}}
\textbf{\bn{পচা}} — \emph{pôchā} — \textquotedblleft{}rotten\textquotedblright{}.

\begin{grammarlens}[title={What you've built}]
One more word for this chapter.
\end{grammarlens}

\subsection*{Your turn}
\begin{itemize}
\raggedright
  \item \emph{Say it:} \emph{pôchā}
  \item \emph{Say it:} \emph{pôchā}, once more
  \item \emph{Say it:} say \emph{pôchā}
  \item \emph{Recall:} say the Bengali for annoyed, then the Bengali for fragrant, then say \emph{pôchā} again
\end{itemize}

\begin{tcolorbox}[breakable,colback=teal!4,colframe=teal!35!black,title={Before you move on}]
What does \bn{পচা} mean? (\textquotedblleft{}Rotten\textquotedblright{}.) Say it once more.
\end{tcolorbox}
\section[\texorpdfstring{\bn{আধা} (ādhā)}{আধা (ādhā)}]{\bn{আধা} (ādhā) — half}

\label{lesson:BN-C162-adha}



\begin{quote}
Before the new one: say the Bengali for fragrant, then the Bengali for rotten.
\end{quote}

\subsection*{You'll want to know: \bn{আধা}}
\textbf{\bn{আধা}} — \emph{ādhā} — \textquotedblleft{}half\textquotedblright{}.

\begin{grammarlens}[title={What you've built}]
One more word for this chapter.
\end{grammarlens}

\subsection*{Your turn}
\begin{itemize}
\raggedright
  \item \emph{Say it:} \emph{ādhā}
  \item \emph{Say it:} \emph{ādhā}, once more
  \item \emph{Say it:} say \emph{ādhā}
  \item \emph{Recall:} say the Bengali for fragrant, then the Bengali for rotten, then say \emph{ādhā} again
\end{itemize}

\begin{tcolorbox}[breakable,colback=teal!4,colframe=teal!35!black,title={Before you move on}]
What does \bn{আধা} mean? (\textquotedblleft{}Half\textquotedblright{}.) Say it once more.
\end{tcolorbox}
\section[\texorpdfstring{\bn{পুরো} (puro)}{পুরো (puro)}]{\bn{পুরো} (puro) — whole, entire}

\label{lesson:BN-C162-puro}



\begin{quote}
Before the new one: say the Bengali for rotten, then the Bengali for half.
\end{quote}

\subsection*{You'll want to know: \bn{পুরো}}
\textbf{\bn{পুরো}} — \emph{puro} — \textquotedblleft{}whole, entire\textquotedblright{}.

\begin{grammarlens}[title={What you've built}]
One more word for this chapter.
\end{grammarlens}

\subsection*{Your turn}
\begin{itemize}
\raggedright
  \item \emph{Say it:} \emph{puro}
  \item \emph{Say it:} \emph{puro}, once more
  \item \emph{Say it:} say \emph{puro}
  \item \emph{Recall:} say the Bengali for rotten, then the Bengali for half, then say \emph{puro} again
\end{itemize}

\begin{tcolorbox}[breakable,colback=teal!4,colframe=teal!35!black,title={Before you move on}]
What does \bn{পুরো} mean? (\textquotedblleft{}Whole, entire\textquotedblright{}.) Say it once more.
\end{tcolorbox}
\section[\texorpdfstring{\bn{সামান্য} (sāmānyô)}{সামান্য (sāmānyô)}]{\bn{সামান্য} (sāmānyô) — slight, a little}

\label{lesson:BN-C162-samanyo}



\begin{quote}
Before the new one: say the Bengali for half, then the Bengali for whole.
\end{quote}

\subsection*{You'll want to know: \bn{সামান্য}}
\textbf{\bn{সামান্য}} — \emph{sāmānyô} — \textquotedblleft{}slight, a little\textquotedblright{}.

\begin{grammarlens}[title={What you've built}]
One more word for this chapter.
\end{grammarlens}

\subsection*{Your turn}
\begin{itemize}
\raggedright
  \item \emph{Say it:} \emph{sāmānyô}
  \item \emph{Say it:} \emph{sāmānyô}, once more
  \item \emph{Say it:} say \emph{sāmānyô}
  \item \emph{Recall:} say the Bengali for half, then the Bengali for whole, then say \emph{sāmānyô} again
\end{itemize}

\begin{tcolorbox}[breakable,colback=teal!4,colframe=teal!35!black,title={Before you move on}]
What does \bn{সামান্য} mean? (\textquotedblleft{}Slight, a little\textquotedblright{}.) Say it once more.
\end{tcolorbox}
